\ifdefined\gkver\else\def\gkver{student}\fi
\documentclass[\gkver]{gaokaozhenti}
\begin{document}

\chapter{2012年重庆理综}
%% paperType: 理综物理部分

\begin{enumerate}
\item
%% number: 14
%% paperName: 2012年重庆理综第 14 题 6 分
%% typeId: 1
%% type: 单选题
%% chapter: 机械振动
%% point: 振动图象
%% method: 无
%% score: 6
%% degree: 650
%% duplicateId: 0
%% body:
装有砂粒的试管竖直静浮于水面，如图所示。将试管竖直提起少许，然后由静止释放并开始计时，在一定时间内试管在竖直方向近似做简谐运动。若取竖直向上为正方向，则以下描述试管振动的图象中可能正确的是\xzanswer{D}
\onepicture[h=2.2cm, align=right]{figs/14.png}
\fourchoices[answer=D, ispicture=true, h=1.8cm]
{figs/14a.png}
{figs/14b.png}
{figs/14c.png}
{figs/14d.png}
%% answer: D
%% memo:
\memoanswer{
根据题中规定的正方向，开始计时时刻位移为正的最大值，由于简谐运动的图象是正弦或余弦曲线，可知 D 正确。
}

\item
%% number: 15
%% paperName: 2012年重庆理综第 15 题 6 分
%% typeId: 1
%% type: 单选题
%% chapter: 交流电
%% point: 变压器
%% method: 无
%% score: 6
%% degree: 650
%% duplicateId: 0
%% body:
如图所示，理想变压器的原线圈接 $u=11000\sqrt 2 \sin100\pi t$（V）的交变电压，副线圈通过电阻 $r=6\,\UO$ 的导线对“220 V/880 W”的电器 $R_{L}$ 供电，该电器正常工作。由此可知\xzanswer{C}
\onepicture[h=2.6cm, align=right]{figs/15.png}
\fourchoices[answer=C]
{原、副线圈的匝数比为$50:1$}
{交变电压的频率为$100\UHz$}
{副线圈中电流的有效值为$4\UA$}
{变压器的输入功率为$880\UW$}
%% answer: C
%% memo:
\memoanswer{
输入电压的有效值为 $11000\UV$，用电器的额定电压为 $220\UV$，所以变压器的输出电压大于 $220\UV$，原、副线圈的匝数比小于 $\frac{50}{1}$，A 错误；由输入电压的表达式知，$f=\frac{100\pi}{2\pi}\UHz=50\UHz$，B 错误；副线圈中的电流与用电器中的电流相同，$I=\frac{P}{U}=4\UA$，C 正确；变压器的输出功率为用电器的功率和导线电阻损耗的功率之和，大于 $880\UW$，所以变压器的输入功率大于 $880\UW$，D 错误。
}

\item
%% number: 16
%% paperName: 2012年重庆理综第 16 题 6 分
%% typeId: 1
%% type: 单选题
%% chapter: 气体性质
%% point: 理想气体状态方程
%% method: 无
%% score: 6
%% degree: 650
%% duplicateId: 0
%% body:
图为伽利略设计的一种测温装置示意图，玻璃管的上端与导热良好的玻璃泡连通，下端插入水中，玻璃泡中封闭有一定量的空气。若玻璃管内水柱上升，则外界大气的变化可能是\xzanswer{A}
\onepicture[h=3.0cm, align=right]{figs/16.png}
\fourchoices[answer=A]
{温度降低，压强增大}
{温度升高，压强不变}
{温度升高，压强减小}
{温度不变，压强减小}
%% answer: A
%% memo:
\memoanswer{
设玻璃泡中气体压强为 $p$，外界大气压强为 $p'$，则 $p'=p+\rho gh$，且玻璃泡中气体与外界大气温度相同。液柱上升，气体体积减小，根据理想气体的状态方程 $\frac{pV}{T}=C$ 可知，$\frac{p}{T}$ 变大，即 $\frac{p'}{T}$ 变大，B、C、D 均不符合要求，A 正确。
}

\item
%% number: 17
%% paperName: 2012年重庆理综第 17 题 6 分
%% typeId: 1
%% type: 单选题
%% chapter: 运动和力
%% point: 动量守恒定律
%% method: 无
%% score: 6
%% degree: 450
%% duplicateId: 0
%% body:
质量为 $m$ 的人站在质量为 2$m$ 的平板小车上，以共同的速度在水平地面上沿直线前行，车所受地面阻力的大小与车对地面压力的大小成正比。当车速为 $v_{0}$ 时，人从车上以相对于地面大小为 $v_{0}$ 的速度水平向后跳下。跳离瞬间地面阻力的冲量忽略不计，则能正确表示车运动的 $v-t$ 图象为\xzanswer{B}
\fourchoices[answer=B, ispicture=true, h=1.8cm]
{figs/17a.png}
{figs/17b.png}
{figs/17c.png}
{figs/17d.png}
%% answer: B
%% memo:
\memoanswer{
人跳车前，人和车以大于 $v_{0}$ 的初速度做匀减速直线运动，加速度大小为 $a=\frac{\mu\times 3mg}{3m}=\mu g$；人跳车瞬间，人和车组成的系统动量守恒，规定初速度方向为正方向，则 $3mv_{0}=-mv_{0}+2mv$，得 $v=2v_{0}$，此后车做减速运动的加速度 $a'=\frac{\mu\times 2mg}{2m}=\mu g=a$，B 项正确。
}

\item
%% number: 18
%% paperName: 2012年重庆理综第 18 题 6 分
%% typeId: 1
%% type: 单选题
%% chapter: 曲线运动
%% point: 双星问题
%% method: 无
%% score: 6
%% degree: 450
%% duplicateId: 0
%% body:
冥王星与其附近的另一星体卡戎可视为双星系统，质量比约为7$:$1，同时绕它们连线上某点O做匀速圆周运动。由此可知，冥王星绕O点运动的\xzanswer{A}
\fourchoices[answer=A]
{轨道半径约为卡戎的$\frac{1}{7}$}
{角速度大小约为卡戎的$\frac{1}{7}$}
{线速度大小约为卡戎的7倍}
{向心力大小约为卡戎的7倍}
%% answer: A
%% memo:
\memoanswer{
双星系统内的两颗星运动的角速度相同，B 错误；双星的向心力为二者间的万有引力，所以向心力大小也相同，D 错误；根据 $m_{1}\omega^{2}r_{1}=m_{2}\omega^{2}r_{2}$，得 $\frac{r_{1}}{r_{2}}=\frac{m_{2}}{m_{1}}=\frac{1}{7}$，A 正确；根据 $v=\omega r$，得 $\frac{v_{1}}{v_{2}}=\frac{r_{1}}{r_{2}}=\frac{1}{7}$，C 错误。
}

\item
%% number: 19
%% paperName: 2012年重庆理综第 19 题 6 分
%% typeId: 1
%% type: 单选题
%% chapter: 原子和原子核
%% point: 人工核反应
%% method: 无
%% score: 6
%% degree: 650
%% duplicateId: 0
%% body:
以下是物理学史上3个著名的核反应方程：\ce{x + ^{7}_{3}Li -> 2y}，\ce{y + ^{14}_{7}N -> x + ^{17}_{8}O}，\ce{y + ^{9}_{4}Be -> z + ^{12}_{6}C}，$x$、$y$和$z$是3种不同的粒子，其中$z$是\xzanswer{C}
\fourchoices[answer=C]
{$\alpha$粒子}
{质子}
{中子}
{电子}
%% answer: C
%% memo:
\memoanswer{
将上述三个方程相加，整理后得 $\ce{^{7}_{3}Li + ^{14}_{7}N + ^{9}_{4}Be -> ^{17}_{8}O + ^{12}_{6}C + z}$，根据电荷数守恒和质量数守恒，$z$ 的质量数为 1，电荷数为 0，为中子，C 正确。
}

\item
%% number: 20
%% paperName: 2012年重庆理综第 20 题 6 分
%% typeId: 1
%% type: 单选题
%% chapter: 电场
%% point: 等势面
%% method: 无
%% score: 6
%% degree: 450
%% duplicateId: 0
%% body:
空间中 P、Q 两点处各固定一个点电荷，其中 P 点处为正电荷，P、Q 两点附近电场的等势面分布如图所示，a、b、c、d 为电场中的 4 个点，则\xzanswer{D}
\onepicture[h=2.6cm, align=right]{figs/20.png}
\fourchoices[answer=D]
{P、Q 两点处的电荷等量同种}
{a 点和 b 点的电场强度相同}
{c 点的电势低于 d 点的电势}
{负电荷从 a 到 c，电势能减少}
%% answer: D
%% memo:
\memoanswer{
由图中等势面的对称性知，P、Q 两处为等量异种电荷，A 错误；由于电场线与等势面垂直，所以 a、b 两处的电场强度方向不同，B 错误；P 处为正电荷，c 在离 P 更近的等势面上，c 点的电势高于 d 点的电势，C 错误；从 a 到 c，电势升高，负电荷电势能减少，D 正确。
}

\item
%% number: 21
%% paperName: 2012年重庆理综第 21 题 6 分
%% typeId: 1
%% type: 单选题
%% chapter: 电磁感应
%% point: 电磁感应中图象
%% method: 无
%% score: 6
%% degree: 450
%% duplicateId: 0
%% body:
如图所示，正方形区域MNPQ内有垂直纸面向里的匀强磁场。在外力作用下，一正方形闭合刚性导线框沿QN方向匀速运动，$t=$0时刻，其四个顶点M$'$、N$'$、P$'$、Q$'$恰好在磁场边界中点。下列图象中能反映线框所受安培力$f$的大小随时间$t$变化规律的是\xzanswer{B}
\onepicture[h=2.6cm, align=right]{figs/21.png}
\fourchoices[answer=B, ispicture=true, h=1.8cm]
{figs/21a.png}
{figs/21b.png}
{figs/21c.png}
{figs/21d.png}
%% answer: B
%% memo:
\memoanswer{
第一段时间从初位置到 $M'N'$ 离开磁场，图甲表示该过程的任意一个位置，切割磁感线的有效长度为 $M_{1}A$ 与 $N_{1}B$ 之和，即为 $M_{1}M'$ 长度的 2 倍，此时电动势 $E=2Bv^{2}t$，线框受的安培力 $f=2BIvt=\frac{4B^{2}v^{3}t^{2}}{R}$，图象是开口向上的抛物线，C、D 错误；如图乙所示，线框的右端 $M_{2}N_{2}$ 刚好出磁场时，左端 $Q_{2}P_{2}$ 恰与MP共线，此后一段时间内有效长度不变，一直到线框的左端与 $M'N'$ 重合，这段时间内电流不变，安培力大小不变；最后一段时间如图丙所示，从匀速运动至 $M_{2}N_{2}$ 开始计时，有效长度为 $A'C'=l-2vt'$，电动势 $E'=B(l-2vt')v$，线框受的安培力 $F'=\frac{B^{2}(l-2vt')^{2}v}{R}$，图象是开口向上的抛物线，A 错误，B 正确。
\onepicture[h=2.6cm, align=right]{figs/21_an1.png}
\onepicture[h=2.6cm, align=right]{figs/21_an2.png}
\onepicture[h=2.6cm, align=right]{figs/21_an3.png}
}

\item
%% number: 22
%% paperName: 2012年重庆理综第 22 题 9 分
%% typeId: 4
%% type: 实验题
%% chapter: 光学
%% point: 测定玻璃折射率
%% method: 无
%% score: 9
%% degree: 650
%% duplicateId: 0
%% body:
如图所示为光学实验用的长方体玻璃砖，它的\_\_\_\_\_\_\_\_面不能用手直接接触。

在用插针法测定玻璃砖折射率的实验中，两位同学绘出的玻璃砖和三个针孔 a、b、c 的位置相同，且插在 c 位置的针正好挡住插在 a、b 位置的针的像，但最后一个针孔的位置不同，分别为 d、e 两点，如图所示，计算折射率时，用\_\_\_\_\_\_\_\_（填“d”或“e”）得到的值较小，用\_\_\_\_\_\_\_\_（填“d”或“e”）点得到的值误差较小。
\twopicture[align=right, hA=0.95cm, hB=2.7cm]
{figs/22a.png}
{figs/22b.png}
%% answer: 
\jdanswer{
光学，d，e
}
%% memo:
\memoanswer{
玻璃砖的光学面不能用手直接接触，接触面的污渍会影响接触面的平整，进而影响折射率的测定。连接 dc、ec 并延长至玻璃砖的光学面与白纸的交线，交点为出射点，入射点与出射点的连线即为折射光线，入射角一定，用 d 点时，折射角大，折射率小；对于两光学面平行的玻璃砖，入射光线和出射光线平行，ec 连线与入射光线平行，误差小。
}

\item
%% number: 22
%% paperName: 2012年重庆理综第 22 题 10 分
%% typeId: 4
%% type: 实验题
%% chapter: 电路
%% point: 测电动势与内阻
%% method: 无
%% score: 10
%% degree: 450
%% duplicateId: 0
%% body:
某中学生课外科技活动小组利用铜片、锌片和家乡盛产的柑橘制作了果汁电池，他们测量这种电池的电动势$E$和内阻$r$，并探究电极间距对$E$和$r$的影响，实验器材如图所示。
\onepicture[h=3.1cm, align=right]{figs/22c.png}
\begin{enumerate}
	\item
	测量$E$和$r$的实验方案为：调节滑动变阻器，改变电源两端的电压$U$和流过电源的电流$I$，依据公式\_\_\_\_\_\_，利用测量数据作出$U-I$图象，得出$E$和$r$。
	\item
	将电压表视为理想表，要求避免电流表分压作用对测量结果的影响，请在图中用笔画线代替导线连接电路。
	\item
	实验中依次减小铜片与锌片的间距，分别得到相应果汁电池的$U-I$图象如图中（a）（b）（c）（d）所示，由此可知：
\end{enumerate}
\onepicture[h=3.4cm, align=right]{figs/22d.png}
在该实验中，随电极间距的减小，电源电动势\_\_\_\_\_\_（填“增大”、“减小”或“不变”），电源内阻\_\_\_\_\_\_（填“增大”、“减小”或“不变”）。曲线（$c$）对应的电源电动势$E=$\_\_\_\_\_\_\_\_\_\_V，内阻$r=$\_\_\_\_\_\_\_\_$\Omega$，当外电路总电阻为2500 $\Omega$时，该电源的输出功率$P=$\_\_\_\_\_\_\_\_mW。（均保留三位有效数字）
%% answer: 
\jdanswer{
\begin{enumerate}
	\item
	$U=E-Ir$

	\item
	如图所示。

	\item
	不变，增大，0.975，478，0.268

\end{enumerate}
}
%% memo:
\memoanswer{
①根据闭合电路的欧姆定律，$U=E-Ir$。

②如图所示。
\onepicture[h=3.1cm, align=right]{figs/22c_an.png}

③图线与纵轴的交点的纵坐标即为电动势，交点纵坐标不变，所以电动势不变，为 $0.975\UV$；图线斜率的绝对值即为内阻，斜率的绝对值变大，内阻变大。（c）对应图线斜率绝对值约为 478（从线上选取两点即可求斜率），电源的输出功率 $P=\left(\frac{E}{R+r}\right)^{2}R=\left(\frac{0.975}{2500+478}\right)^{2}\times 2500\,\UW=0.268\,\UmW$。
}

\item
%% number: 23
%% paperName: 2012年重庆理综第 23 题 16 分
%% typeId: 6
%% type: 计算题
%% chapter: 功和能
%% point: 动能定理
%% method: 无
%% score: 16
%% degree: 450
%% duplicateId: 0
%% body:
如图所示为一种摆式摩擦因数测量仪，可测量轮胎与地面间动摩擦因数，其主要部件有：底部固定有轮胎橡胶片的摆锤和连接摆锤的轻质细杆。摆锤的质量为$m$，细杆可绕轴$O$在竖直平面内自由转动，摆锤重心到$O$点距离为$L$。测量时，测量仪固定于水平地面，将摆锤从与$O$等高的位置处静止释放。摆锤到最低点附近时，橡胶片紧压地面擦过一小段距离$s$（$s<L$），之后继续摆至与竖直方向成$\theta$角的最高位置。若摆锤对地面的压力可视为大小为$F$的恒力，重力加速度为$g$，求：
\begin{enumerate}
	\item
	摆锤在上述过程中损失的机械能；
	\item
	在上述过程中摩擦力对摆锤所做的功；
	\item
	橡胶片与地面之间的动摩擦因数。
\end{enumerate}
\onepicture[h=3.4cm, align=right]{figs/23.png}
%% answer: 
\jdanswer{
\begin{enumerate}
	\item
	$mgL\cos\theta$

	\item
	$-mgL\cos\theta$

	\item
	$\frac{mgL\cos \theta }{Fs}$

\end{enumerate}
}
%% memo:
\memoanswer{
（1）机械能的损失：$\Delta E=E_{1}-E_{2}$

损失的机械能：$\Delta E=E_{1}-E_{2}=mgL-mgL(1-\cos\theta)=mgL\cos\theta$

（2）由动能定理得：$mgL\cos\theta+W_{f}=0$

摩擦力对摆锤所做的功 $W_{f}=-mgL\cos \theta$

（3）由于 $W_{f}=-\mu Fs$

橡胶片与地面之间的动摩擦因数：$\mu=\frac{mgL\cos\theta}{Fs}$
}

\item
%% number: 24
%% paperName: 2012年重庆理综第 24 题 18 分
%% typeId: 6
%% type: 计算题
%% chapter: 磁场
%% point: 带电粒子在复合场中的运动
%% method: 无
%% score: 18
%% degree: 450
%% duplicateId: 0
%% body:
有人设计了一种带电颗粒的速率分选装置，其原理如图所示，两带电金属板间有匀强电场，方向竖直向上，其中$PQNM$矩形区域内还有方向垂直纸面向外的匀强磁场。一束比荷（电荷量与质量之比）均为$\frac{1}{k}$的带正电颗粒，以不同的速率沿着磁场区域的水平中心线为$O'O$进入两金属板之间，其中速率为$v_{0}$的颗粒刚好从$Q$点处离开磁场，然后做匀速直线运动到达收集板。重力加速度为$g$，$PQ=3d$，$NQ=2d$，收集板与$NQ$的距离为$l$，不计颗粒间相互作用。求：
\begin{enumerate}
	\item
	电场强度$E$的大小；
	\item
	磁感应强度$B$的大小；
	\item
	速率为$\lambda v_{0}$（$\lambda>1$）的颗粒打在收集板上的位置到$O$点的距离。
\end{enumerate}
\onepicture[h=3.4cm, align=right]{figs/24.png}
%% answer: 
\jdanswer{
\begin{enumerate}
	\item
	$kg$

	\item
	$\frac{kv_{0}}{5d}$

	\item
	$d(5\lambda  - \sqrt {25\lambda ^{2} - 9} ) + \frac{3l}{\sqrt {25\lambda ^{2} - 9} }$

\end{enumerate}
}
%% memo:
\memoanswer{
（1）设带电颗粒的电荷量为 $q$，质量为 $m$。有

$Eq=mg$

将 $\frac{q}{m}=\frac{1}{k}$ 代入，得 $E=kg$

（2）如图，有

$qv_{0}B=m\frac{v_{0}^{2}}{R}$

$R^{2}=(3d)^{2}+(R-d)^{2}$

得 $B=\frac{kv_{0}}{5d}$

\onepicture[h=3.2cm, align=right]{figs/24_an1.png}

（3）如图所示，有

$q\lambda v_{0}B=m\frac{(\lambda v_{0})^{2}}{R_{1}}$

$\tan\theta=\frac{3d}{\sqrt{R_{1}^{2}-(3d)^{2}}}$

$y_{1}=R_{1}-\sqrt{R_{1}^{2}-(3d)^{2}}$

$y_{2}=l\tan\theta$

$y=y_{1}+y_{2}$

得 $y=d(5\lambda-\sqrt{25\lambda^{2}-9})+\frac{3l}{\sqrt{25\lambda^{2}-9}}$

\onepicture[h=3.2cm, align=right]{figs/24_an2.png}
}

\item
%% number: 25
%% paperName: 2012年重庆理综第 25 题 19 分
%% typeId: 6
%% type: 计算题
%% chapter: 运动和力
%% point: 牛顿定律的应用
%% method: 无
%% score: 19
%% degree: 250
%% duplicateId: 0
%% body:
某校举行托乒乓球跑步比赛，赛道为水平直道，比赛距离为 $s$。比赛时，某同学将球置于球拍中心，以大小为 $a$ 的加速度从静止开始做匀加速直线运动，当速度达到 $v_{0}$ 时，再以 $v_{0}$ 做匀速直线运动跑至终点。整个过程中球一直保持在球拍中心不动。比赛中，该同学在匀速直线运动阶段保持球拍的倾角为 $\theta_{0}$，如图所示，设球在运动中受到的空气阻力大小与其速度大小成正比，方向与运动方向相反，不计球与球拍之间的摩擦，球的质量为 $m$，重力加速度为 $g$。
\begin{enumerate}
	\item
	求空气阻力大小与球速大小的比例系数 $k$；
	\item
	求在加速跑阶段球拍倾角 $\theta$ 随速度 $v$ 变化的关系式；
	\item
	整个匀速跑阶段，若该同学速度仍为 $v_{0}$，而球拍的倾角比 $\theta_{0}$ 大了 $\beta$ 并保持不变，不计球在球拍上的移动引起的空气阻力变化，为保证到达终点前球不从球拍上距离中心为 $r$ 的下边沿掉落，求 $\beta$ 应满足的条件。
\end{enumerate}
\onepicture[h=2.6cm, align=right]{\includesvg{figs/25}}
%% answer: 
\jdanswer{
\begin{enumerate}
	\item
	$\frac{mg\tan {\theta _0}}{v_{0}}$

	\item
	$\tan\theta = \frac{a}{g} + \frac{v}{v_{0}}\tan\theta_{0}$

	\item
	$\sin\beta \leq \frac{2r\cos {\theta _0}}{g{{\left(\frac{s}{v_{0}} - \frac{v_{0}}{2a}\right)}^2}}$

\end{enumerate}
}
%% memo:
\memoanswer{
（1）在匀速运动阶段，有 $mg\tan\theta_{0}=kv_{0}$

得 $k=\frac{mg\tan\theta_{0}}{v_{0}}$

（2）加速阶段，设球拍对球的支持力为 $N'$，有

$N'\sin\theta-kv=ma$

$N'\cos\theta=mg$

得 $\tan\theta=\frac{a}{g}+\frac{v}{v_{0}}\tan\theta_{0}$

（3）以速度 $v_{0}$ 匀速运动时，设空气阻力与重力的合力为 $F$，有 $F=\frac{mg}{\cos\theta_{0}}$

球拍倾角为 $\theta_{0}+\beta$ 时，空气阻力与重力的合力不变，设球沿球拍面下滑的加速度大小为 $a'$，有

$F\sin \beta=ma'$

设匀速跑阶段所用时间为 $t$，有

$t=\frac{s}{v_{0}}-\frac{v_{0}}{2a}$

球不从球拍上掉落的条件 $\frac{1}{2}a't^{2}\leq r$

得 $\sin \beta \leq \frac{2r\cos \theta_{0}}{g\left(\frac{s}{v_{0}}-\frac{v_{0}}{2a}\right)^{2}}$
}

\end{enumerate}

\end{document}
